\documentclass{article}
\usepackage{amsmath, amssymb, amsthm}
\usepackage{hyperref}
\usepackage{geometry}
\geometry{margin=1in}

\title{Erd\H{o}s Problem \#217}
\date{Last updated: 2025-08-31}
\author{Source: erdosproblems.com}

\begin{document}
\maketitle

\noindent\textbf{Status:} OPEN \quad \textbf{No prize}

\noindent\textbf{Tags:} geometry, distances

\medskip

\noindent\textbf{Problem Statement:}

For which $n$ are there $n$ points in $\mathbb{R}^2$, no three on a line and no four on a circle, which determine $n-1$ distinct distances and so that (in some ordering of the distances) the $i$th distance occurs $i$ times?




An example with $n=4$ is an isosceles triangle with the point in the centre. Erd\H{o}s originally believed this was impossible for $n\geq 5$, but Pomerance constructed a set with $n=5$ (see \cite{Er83c} for a description). Pal\'{a}sti \cite{Pa87} constructed such a set with $n=7$. According to \cite{Pa89} when Erd\H{o}s saw this he asked whether there exists such a set of $5$ points with no equilateral triangles; Pal\'{a}sti \cite{Pa89} {IMAGE=217a,constructed a set} of $n=6$ points, satisfying these restrictions, with no equilateral triangles. Pal\'{a}sti \cite{Pa89b} has {IMAGE=217b,constructed a set} with $n=8$. Erd\H{o}s believed this is impossible for all sufficiently large $n$. This would follow from $h(n)\geq n$ for sufficiently large $n$, where $h(n)$ is as in [98] .




References

[Er83c] Erd\H{o}s, Paul, Combinatorial problems in geometry . Math. Chronicle (1983), 35-54.

[Pa87] Pal\'{a}sti, Ilona, On the seven points problem of {P}. {E}rd\H os . Studia Sci. Math. Hungar. (1987), 447--448.

[Pa89] Pal\'{a}sti, Ilona, A distance problem of {P}. {E}rd\H os with some further restrictions . Discrete Math. (1989), 155--156.

[Pa89b] Pal\'{a}sti, I., Lattice-point examples for a question of {E}rd\H os . Period. Math. Hungar. (1989), 231--235.

\noindent\textbf{Related OEIS sequences:} possible


\bigskip
\noindent\small{Source: \url{https://www.erdosproblems.com/217}}
\end{document}
